% -------------------------------------------------------------------------
% WBS ID: RES-GEO-CST
% Deliverable: Geometric Feasibility and Design Constraints
% Status: In progress
% Evidence status: Fixed-rigid geometry feasibility gates established
% Review status: Not reviewed
% Last updated: 2026-07-19
% Dependencies: RES-GEO-PAR, RES-DAT-GEO, RES-NUM-MSH
% Downstream interfaces: RES-GEO-SPEC, RES-NUM-MSH, Software geometry manifest
% -------------------------------------------------------------------------

\section{RES-GEO-CST --- Geometric Feasibility and Design Constraints}
\label{sec:res-geo-cst}

\subsection{Purpose and Scope}
\label{sec:res-geo-cst-purpose}

This Deliverable narrows admissible geometry cases to fixed-rigid
barrier-class comparisons that can be represented on the accepted
terrain, meshed robustly and interpreted through hydrodynamic outputs.
Material limits, foundation constraints, deformability and damage are
recorded as stretch structural concerns rather than active geometric
feasibility criteria.

\subsection{Research Questions}
\label{sec:res-geo-cst-research-questions}

% TODO: State the questions that must be answered before the Deliverable
% can be accepted.

\subsection{Terminology and Definitions}
\label{sec:res-geo-cst-definitions}

% TODO: Define the technical terminology, symbols, quantities and
% classifications used in this Deliverable.

\subsection{Literature Evidence}
\label{sec:res-geo-cst-literature-evidence}

% TODO: Record evidence from peer-reviewed papers, authoritative datasets,
% standards, technical reports and primary documentation.
%
% Do not create a detached literature-summary list. Explain how each source
% supports, contradicts or limits a project decision.

\subsection{Governing Relations and Quantitative Definitions}
\label{sec:res-geo-cst-governing-relations}

% TODO: Record relevant equations, dimensionless groups, empirical models,
% extraction rules or quantitative definitions.
%
% Use align environments for displayed equations.
% Every displayed equation must have a unique \label{eq:...}.
% Do not place punctuation after displayed equations.

\subsection{Geometry Family or Representation}
\label{sec:res-geo-cst-geometry-family-representation}

% TODO: Complete this RES-GEO Domain-specific subsection.

\subsection{Design Variables and Parameter Bounds}
\label{sec:res-geo-cst-design-variables-parameter-bounds}

% TODO: Complete this RES-GEO Domain-specific subsection.

\subsection{Geometric Descriptors}
\label{sec:res-geo-cst-geometric-descriptors}

% TODO: Complete this RES-GEO Domain-specific subsection.

\subsection{Placement and Arrangement Rules}
\label{sec:res-geo-cst-placement-arrangement-rules}

% TODO: Complete this RES-GEO Domain-specific subsection.

\subsection{Feasibility Constraints}
\label{sec:res-geo-cst-feasibility-constraints}

Candidate interventions shall be rejected or marked unresolved if they cannot
be constructed plausibly, cannot be represented in the selected numerical
workflow, violate known material or foundation constraints, or lack enough
provenance to support simulation.

For the active hydrodynamic baseline, a geometry shall also be rejected
or deferred if it requires unresolved porous resistance, moving
structural boundaries, scour-induced terrain change or material
damage to explain its intended performance.

\textbf{Selected approach.} Feasibility is first a data and numerical
representation gate: the case must fit within the accepted corridor
polygon, avoid sponge zones, preserve mesh quality near walls and have
a traceable descriptor record. Supporting evidence:
\textcite{watanabe2022validation} for validation-facing geometry
control and \textcite{deljesus2012porous} for separating impermeable
and porous/dissipating assumptions.

\subsection{Two- and Three-Dimensional Representation}
\label{sec:res-geo-cst-two-three-dimensional-representation}

% TODO: Complete this RES-GEO Domain-specific subsection.

\subsection{Candidate Approaches}
\label{sec:res-geo-cst-candidate-approaches}

% TODO: Define the credible candidate approaches considered.

\subsection{Critical Comparison}
\label{sec:res-geo-cst-critical-comparison}

% TODO: Compare candidates using physically and project-relevant criteria.
% Do not select an approach solely on popularity or implementation ease.

\subsection{Assumptions and Applicability}
\label{sec:res-geo-cst-assumptions}

% TODO: State assumptions and the conditions under which they are valid.

\subsection{Limitations and Failure Conditions}
\label{sec:res-geo-cst-limitations}

% TODO: State where the evidence, model or selected approach ceases to be
% reliable or applicable.

\subsection{Project-Specific Conclusions}
\label{sec:res-geo-cst-conclusions}

The Studentship shall not spend simulation effort on barrier shapes
whose geometry cannot be reproduced, clipped to the corridor and
post-processed for hydrodynamic metrics.

\subsection{Selected Approach and Rationale}
\label{sec:res-geo-cst-selected-approach}

% TODO: Record the selected approach only after sufficient evidence exists.
% State the alternatives rejected and the reasons for rejection.

\subsection{Unresolved Decisions}
\label{sec:res-geo-cst-unresolved-decisions}

Open decisions are final mesh-quality limits, minimum wall clearance
from artificial boundaries, and whether any dissipating concept is
admissible without a porous closure.

\subsection{Requirements and Handoffs}
\label{sec:res-geo-cst-requirements}

Software shall report a geometry-feasibility status before launching
local runs, including corridor intersection checks, minimum cell-size
estimates, wall clearance and unresolved-physics flags.

\subsection{Deliverable Completion Record}
\label{sec:res-geo-cst-completion-record}

% TODO: Record whether the Deliverable is:
%
% - Not started
% - In progress
% - Draft complete
% - Reviewed
% - Accepted
%
% Acceptance requires:
%
% - the research questions to be answered;
% - claims to be supported by appropriate evidence;
% - assumptions and limitations to be explicit;
% - the project-specific conclusion to be stated;
% - downstream requirements to be traceable;
% - unresolved decisions to be closed or formally deferred.
